```latex
\section{Evaluation Methods for Code Translation}

This section focuses on evaluation metrics used for code translation, rather
than benchmark datasets themselves. While benchmarks define the evaluation
setup, metrics determine how the quality of generated code is measured.

Evaluation metrics for code translation can be broadly divided into
similarity-based and execution-based approaches.

Text-based metrics are derived from natural language processing and measure the
similarity between generated and reference code at the token level. One of the
most widely used metrics is BLEU, which evaluates n-gram overlap between the
generated output and the reference translation \cite{bleu_metric}. BLEU is easy
to compute and remains widely used in both natural language translation and
code generation tasks. However, it does not account for syntactic structure,
semantic equivalence, or executability. Different implementations of the same
logic may receive very different BLEU scores, while code with high lexical
similarity may still fail to compile or execute correctly
\cite{tran2019ruby, evtikhiev2022bleu}.

To address these limitations, CodeBLEU was introduced as an extension of BLEU
specifically designed for code evaluation. In addition to n-gram overlap,
CodeBLEU incorporates abstract syntax tree (AST) structure and data-flow
information, providing a more comprehensive assessment of code quality
\cite{ren2020codebleu}. Ren et al. showed that CodeBLEU correlates better with
programmer judgments than BLEU across multiple code-related tasks, including
code translation, code refinement, and text-to-code generation
\cite{ren2020codebleu}. However, CodeBLEU still depends on similarity to a
reference solution and may fail to capture whether the generated code is
actually functionally correct.

Other similarity-based metrics have also been proposed. For example,
CrystalBLEU attempts to reduce the influence of frequently occurring code
fragments, while CodeBERTScore uses learned code representations to estimate
semantic similarity between generated and reference code. More recent work has
also explored metrics such as CodeScore and ICE-Score, which use large language
models to estimate the quality and correctness of generated code
\cite{codescore2023, ice_score}. Although these methods often correlate better
with human judgments than BLEU, they still do not fully replace direct
execution-based testing.

Execution-based evaluation methods assess whether the generated code behaves
correctly by compiling and running it against predefined test cases. Metrics
such as pass rate or pass@k directly measure functional correctness. These
methods are particularly important for code translation because the main goal is
not only to produce code that resembles the reference output, but also to
preserve the original program behavior.

Several recent studies have shown that models achieving high scores on
similarity-based metrics do not necessarily perform well in execution-based
evaluation. Huang et al. demonstrated that models with strong BLEU and
CodeBLEU performance often fail when evaluated using execution-based metrics,
highlighting the gap between surface-level similarity and true functional
correctness \cite{huang2022exeds}. Similar findings have also been reported in
other studies, which show that lexical overlap often correlates poorly with
human judgments and program behavior \cite{evtikhiev2022bleu,
tran2019ruby}.

Despite their advantages, execution-based methods also have limitations. Their
effectiveness depends heavily on the quality and coverage of the available test
cases. A program may pass all tests while still containing hidden logical
errors if the test suite is incomplete. In addition, execution-based evaluation
is significantly more computationally expensive than similarity-based metrics,
because it requires compilation, sandboxing, and repeated program execution
\cite{huang2022exeds}. Existing execution benchmarks also typically focus on
isolated functions or short code snippets rather than realistic multi-file
software systems, which limits their ability to represent real-world code
migration scenarios \cite{evaluation_limitations, llm_limitations}.

These limitations indicate that there is still no universally reliable
evaluation method for code translation. Similarity-based metrics are efficient
but often fail to capture semantic correctness, while execution-based methods
provide more realistic evaluation but remain expensive and dependent on the
quality of the benchmark and test cases. This motivates the use of both types
of evaluation in this thesis.
